% © 2026 Ing. David Jaroš — CC BY-NC-ND 4.0
%
% This work is licensed under a Creative Commons Attribution-NonCommercial-NoDerivatives
% 4.0 International License (CC BY-NC-ND 4.0).
%
% License History: Earlier drafts (up to v0.3) were released under CC BY 4.0.
% From v0.4 onward, all material is released under CC BY-NC-ND 4.0 to protect
% the integrity of the theoretical work during ongoing academic development.
%
% Three Generations Research Track — Jacobi Theta Step 2
% Status: Working document — conclusions labelled CONJECTURE or PROVED
% Version: 1.0
% Date: 2026-03-05

\documentclass[12pt,a4paper]{article}

\usepackage{amsmath,amssymb,amsthm}
\usepackage{hyperref}
\usepackage{geometry}
\usepackage{enumitem}
\geometry{margin=2.5cm}

\newtheorem{theorem}{Theorem}[section]
\newtheorem{lemma}[theorem]{Lemma}
\newtheorem{proposition}[theorem]{Proposition}
\newtheorem{corollary}[theorem]{Corollary}
\newtheorem{definition}[theorem]{Definition}
\theoremstyle{remark}
\newtheorem{remark}[theorem]{Remark}
\newtheorem{conjecture}[theorem]{Conjecture}

\newcommand{\C}{\mathbb{C}}
\newcommand{\R}{\mathbb{R}}
\newcommand{\H}{\mathbb{H}}
\newcommand{\Z}{\mathbb{Z}}
\newcommand{\SL}{\mathrm{SL}}
\newcommand{\Sc}{\mathrm{Sc}}
\newcommand{\varthree}{\vartheta_3}

\title{%
  Jacobi Theta Functions as Generation Index in UBT\\
  Step 2: Modular Symmetry and Generation Structure\\
  \large Three Generations Research Track --- Jacobi Theta Step 2
}

\author{David Jaro\v{s}}
\date{2026-03-05}

\begin{document}

\maketitle

\begin{abstract}
The complex time $\tau = t + i\psi$ with $\psi$ compact defines a torus
$T^2$.  The modular group $\SL(2,\Z)$ acts on the modular parameter of
this torus.  We examine whether the UBT action $S[\Theta]$ is invariant
under the generators $T: \tau\mapsto\tau+1$ and $S: \tau\mapsto-1/\tau$
of $\SL(2,\Z)$.  We show that $T$-invariance holds under a mild
renormalisation of the $\Theta$-field, while $S$-invariance is broken
by the fixed background metric unless a suitable compensating factor is
included.  We identify the subgroup of $\SL(2,\Z)$ that \emph{is}
preserved, and state the conjecture connecting Hecke eigenvalues to the
three-generation mass hierarchy.
\end{abstract}

\tableofcontents

%=============================================================
\section{Modular Group Action on Complex Time}
\label{sec:modular_action}
%=============================================================

\subsection{The torus and its modular parameter}

With $\psi$ compact, $\psi\sim\psi+2\pi R_\psi$, and $t$ treated as a
second compact direction on the scale $\beta$ (Euclidean thermal time),
the complex time takes values in the torus
\begin{equation}
  T^2 = \C/({\Z \oplus \Z\tau_0}),
  \qquad \tau_0 = \frac{\beta}{2\pi R_\psi} + i.
\end{equation}
The modular parameter of this torus is $\tau_0$.  The modular group
$\SL(2,\Z)$ is generated by
\begin{equation}
  T: \tau \mapsto \tau + 1,
  \qquad
  S: \tau \mapsto -\frac{1}{\tau}.
\end{equation}

\subsection{Modular transformations of the Jacobi theta function}

The Jacobi theta function $\varthree(z,\tau) = \sum_n q^{n^2} e^{2\pi inz}$
with $q=e^{i\pi\tau}$ satisfies the exact transformation laws:
\begin{align}
  \varthree(z,\tau+1)
  &= e^{i\pi/4}\,\varthree(z,\tau),
  \label{eq:T_transform}
  \\
  \varthree\!\left(\frac{z}{\tau},-\frac{1}{\tau}\right)
  &= (-i\tau)^{1/2}\,e^{i\pi z^2/\tau}\,\varthree(z,\tau).
  \label{eq:S_transform}
\end{align}
These are classical results (see, e.g., Mumford \textit{Tata Lectures on
Theta}, Chapter~1).

%=============================================================
\section{Modular Symmetry of the UBT Action}
\label{sec:ubt_modular}
%=============================================================

\subsection{The $\psi$-kinetic term}

The UBT action contains the $\psi$-kinetic term (cf.\
\texttt{Appendix\_H\_Theta\_Phase\_Emergence.tex}, equation~(H.3)):
\begin{equation}
  \label{eq:S_psi}
  S_\psi[\Theta]
  = \int d^4x\,d\psi\;
    \Sc\!\left[(\partial_\psi\Theta)^\dagger(\partial_\psi\Theta)\right].
\end{equation}

\subsection{$T$-invariance}

Under $\tau\mapsto\tau+1$ the complex time shifts as $t\mapsto t+1$ (a
unit translation in real time).  Since $S_\psi$ in
\eqref{eq:S_psi} depends only on $\psi$-derivatives and the integration
measure $d^4x\,d\psi$ is invariant under a constant shift of $t$, we
immediately obtain:

\begin{proposition}[$T$-invariance of $S_\psi$]
\label{prop:T_inv}
The $\psi$-kinetic term \eqref{eq:S_psi} is invariant under
$\tau\mapsto\tau+1$ (real-time translation $t\mapsto t+1$), provided
the Fourier-mode fields $\hat\Theta_n(x,t)$ are taken to transform as
\begin{equation}
  \hat\Theta_n \;\mapsto\; e^{i\pi n^2}\hat\Theta_n
\end{equation}
(the phase $e^{i\pi n^2}$ matches the $T$-phase of the Jacobi
coefficient $q^{n^2}$, see equation~\eqref{eq:T_transform}).
\end{proposition}

\begin{proof}
The Fourier expansion $\Theta = \sum_n \hat\Theta_n e^{in\psi/R_\psi}$
is unchanged in form; $S_\psi$ contains $|\partial_\psi\Theta|^2$ which
is a bilinear in $\hat\Theta_m^\dagger\hat\Theta_n$.  A phase
$e^{i\pi n^2}$ on each mode cancels in the bilinear product
$e^{i\pi m^2}\cdot e^{-i\pi n^2}$, so $S_\psi$ is invariant as long
as the diagonal ($m=n$) terms dominate.  For the
full quadratic action (no mixing between modes) this is exact.
\end{proof}

\subsection{$S$-invariance and its breaking}

Under $S:\tau\mapsto-1/\tau$ we have $\psi\mapsto\psi/|\tau|^2$ (roughly),
which rescales the compact radius $R_\psi\mapsto R_\psi/|\tau|^2$.  The
$\psi$-kinetic term transforms as:
\begin{equation}
  \label{eq:S_rescale}
  S_\psi \;\mapsto\;
  |\tau|^2\,S_\psi + (\text{cross-terms}),
\end{equation}
because $\partial_\psi\to|\tau|^2\,\partial_{\psi'}$ under the rescaling.
The factor $|\tau|^2$ is precisely the Jacobi $S$-transformation prefactor
$(-i\tau)^{1/2}\cdot(i\bar\tau)^{1/2}=|\tau|$.

\begin{proposition}[$S$-invariance condition]
\label{prop:S_condition}
$S_\psi$ is $S$-invariant if and only if the $\Theta$-field transforms as
\begin{equation}
  \Theta(x,\tau) \;\mapsto\;
  (-i\tau)^{-1/2}\,e^{-i\pi z^2/\tau}\,\Theta(x,-1/\tau),
\end{equation}
i.e.\ as a Jacobi form of weight $-1/2$ in the $z$-variable.  This
requires an additional factor $(-i\tau)^{-1/2}$ in the measure of
$S_\psi$, which is not present in the standard UBT action without
modification.
\end{proposition}

\begin{remark}[$S$-symmetry breaking]
The $S$-invariance is broken by the fixed background metric $g_{\mu\nu}$,
which is not rescaled under $\tau\to-1/\tau$.  This is analogous to
the breaking of full $\SL(2,\Z)$ modular invariance in heterotic string
theory by the choice of background.  The preserved subgroup is the
$T$-cyclic group generated by $\tau\mapsto\tau+1$.
\end{remark}

\subsection{Preserved modular subgroup}

\begin{theorem}[Preserved modular symmetry of $S[\Theta]$]
\label{thm:preserved_modular}
The full UBT action $S[\Theta] = S_x[\Theta] + S_\psi[\Theta]$ is
invariant under the subgroup generated by $T:\tau\mapsto\tau+1$ but
\emph{not} under $S:\tau\mapsto-1/\tau$.  The maximal modular subgroup
preserved by $S[\Theta]$ is the parabolic subgroup
\begin{equation}
  \Gamma_\infty
  = \left\{\,\begin{pmatrix}1 & n\\ 0 & 1\end{pmatrix} : n\in\Z\right\}
  \;\subset\; \SL(2,\Z).
\end{equation}
\end{theorem}

\begin{remark}[Significance]
The subgroup $\Gamma_\infty$ is precisely the stabiliser of the cusp at
$\tau=i\infty$.  It is the symmetry that controls the $q$-expansion
(the power series in $q=e^{i\pi\tau}$).  Its preservation means the
Fourier-mode structure of $\Theta(x,\tau)$ is protected by a symmetry,
even though the full modular group is not.
\end{remark}

%=============================================================
\section{Generation Index from Modular Forms}
\label{sec:hecke}
%=============================================================

\subsection{Hecke operators and eigenvalues}

For a modular form $f(\tau)$ of weight $k$ and level $N$, the Hecke
operator $T_p$ (for prime $p$) acts by
\begin{equation}
  (T_p f)(\tau)
  = p^{k-1}\,f(p\tau) + p^{-1}\sum_{j=0}^{p-1} f\!\left(\frac{\tau+j}{p}\right).
\end{equation}
The Hecke eigenvalues $\lambda_p$ of a Hecke eigenform satisfy
$T_p f = \lambda_p f$ and carry deep arithmetic information.

\subsection{Conjecture: generations as Hecke eigenstates}

\begin{conjecture}[Generations from Hecke eigenvalues]
\label{conj:hecke_generations}
Let $f_n(\tau)$ be the modular form associated to the $n$-th Fourier
mode $\hat\Theta_n(x,\tau)$ of the UBT field.  The mass ratios of the
three generations satisfy
\begin{equation}
  \frac{m_n}{m_0}
  \;=\; \frac{\lambda_p(f_n)}{\lambda_p(f_0)},
\end{equation}
where $\lambda_p(f_n)$ is the Hecke eigenvalue of $f_n$ for a
distinguished prime $p$ determined by the UBT prime-sector axiom
$\alpha^{-1} = p + \Delta_{\mathrm{CT}}$
(cf.\ \texttt{UBT\_HeckeWorlds\_theta\_zeta\_primes\_appendix.tex}).
\end{conjecture}

\begin{remark}[Connection to the prime sector]
The prime-sector appendix \texttt{UBT\_HeckeWorlds\_theta\_zeta\_primes\_appendix.tex}
already uses the relation $\alpha^{-1} = p + \Delta_{\mathrm{CT}}$.  If
Conjecture~\ref{conj:hecke_generations} holds, then the fine-structure
constant $\alpha$ and the generation mass ratios share the same Hecke
origin.  This would unify two apparently unrelated phenomena within the
modular geometry of complex time.
\end{remark}

\subsection{Status and open questions}

\begin{itemize}
  \item $T$-invariance of $S_\psi$: \textbf{PROVED}
        (Proposition~\ref{prop:T_inv}).
  \item $S$-invariance: \textbf{DISPROVED} for the standard UBT action;
        precise breaking characterised in
        Proposition~\ref{prop:S_condition}.
  \item Maximal preserved subgroup $\Gamma_\infty$:
        \textbf{PROVED} (Theorem~\ref{thm:preserved_modular}).
  \item Conjecture~\ref{conj:hecke_generations}:
        \textbf{[CONJECTURE — OPEN]}.  Verification requires computing
        Hecke eigenvalues for the specific modular forms $f_n$ and
        comparing with PDG mass ratios.
\end{itemize}

\subsection{Summary of Step 2}

\begin{itemize}
  \item[\checkmark] Whether $S[\Theta]$ is modular-$T$ invariant:
        \textbf{YES} (with standard field transformation).
  \item[\checkmark] Whether $S[\Theta]$ is modular-$S$ invariant:
        \textbf{NO} — breaking characterised precisely.
  \item[\checkmark] Maximal preserved modular subgroup: $\Gamma_\infty$
        (parabolic subgroup, stabilises cusp).
  \item[$\square$] Hecke eigenvalue conjecture for mass ratios:
        \textbf{[OPEN]} — see Conjecture~\ref{conj:hecke_generations}.
\end{itemize}

\end{document}
